\begin{lemma}
For every $\eta>0$ and every positive integer $k$, there is $D_0$ such that,
whenever $D\ge D_0$, $\mathcal H$ is a $k$-uniform $k$-simple hypergraph
with maximum degree at most $D$, and $e\in E(\mathcal H)$, one has
\[
b_{L(\mathcal H),kD}(e)
\le
1-\frac{k-1}{2k^2}
+\frac{(k-1)(k-2)}{6k^3\sqrt{k}}+\eta .
\]
\end{lemma}
\begin{proof}
Fix $\eta>0$ and $k\in\mathbb Z^+$.  We first make the finite local-pattern
reduction required for a universal statement.  For fixed $D$ and an
arbitrary admissible pair $(\mathcal H,e)$, retain $e$, every edge meeting
$e$, and all vertices lying in those retained edges.  The line graph
\noderef{LineGraphOfHypergraph} has edge indices as vertices and adjacency
given by nonempty intersection of the corresponding hyperedges.  Hence the
induced line-graph neighborhood of $e$ is unchanged by this truncation.
Consequently the line-graph degree of $e$, the independent-pair count
\noderef{IndependentPairCount}, the independent-triple count
\noderef{IndependentTripleCount}, and the value of
\noderef{LocalBParameter} at $e$ are unchanged.

The truncation preserves membership in $H(k,k,D)$.  Each retained edge keeps
its original size, so \noderef{UniformHypergraph} is preserved.  Pairwise
edge intersections are original intersections, so \noderef{TSimpleHypergraph}
is preserved.  Vertex degrees only decrease when edge indices are discarded;
using \noderef{HypergraphDegree} and \noderef{MaxDegreeAtMost}, the maximum
degree is still at most $D$.  Thus \noderef{HypergraphClass} still applies.
There are at most $kD$ retained neighboring edge indices and finitely many
retained vertices for fixed $D$ and $k$, so there are only finitely many
truncated local patterns up to relabeling.  Choose a representative
$(\mathcal F,f)$ maximizing $b_{L(\mathcal H),kD}(e)$ among these finitely
many patterns.  Proving the asserted bound for this representative proves it
for the original pair, because the original pair has the same local
parameter as its truncation and the representative is maximal among all
truncated local patterns.
Moreover, every admissible pair has one of these truncations with the same
local parameter, so this representative also satisfies the maximizer
hypothesis appearing in the extremal nodes cited below.

We next split off the case $k=1$.  In a $1$-uniform hypergraph every edge is
a single vertex.  If a neighboring edge $g$ of $e$ meets $e$, then both are
the same singleton as subsets of the vertex set; distinct edge indices may
be parallel, but every two such neighbors meet in that same vertex.  Thus the
line-graph neighborhood of $e$ is a clique.  Consequently there are no
independent pairs and no independent triples in this neighborhood, so
\noderef{IndependentPairCount} and \noderef{IndependentTripleCount} give
$P_{L(\mathcal H)}(e)=T_{L(\mathcal H)}(e)=0$.  The maximum-degree condition
in \noderef{HypergraphClass}, interpreted through
\noderef{HypergraphDegree} and \noderef{MaxDegreeAtMost}, says that the
single vertex of $e$ lies in at most $D$ edge indices; excluding $e$ itself,
\[
\deg_{L(\mathcal H)}(e)\le D-1.
\]
Since in this case the scale in \noderef{LocalBParameter} is $kD=D$, we have
\[
b_{L(\mathcal H),D}(e)
=\frac{\deg_{L(\mathcal H)}(e)}D
\le \frac{D-1}{D}\le 1\le 1+\eta .
\]
For $k=1$ the displayed bound in the statement is exactly $1+\eta$, because
both numerator factors $k-1$ in the correction terms vanish.  This proves
the lemma when $k=1$.

Assume from now on that $k\ge2$.  This is the only case in which we use the
$k$-partite triangle bound, so the denominator $\binom{k}{2}$ is positive.
By \noderef{KUniformKSimpleExtremalSaturation}, after increasing $D_0$ if
necessary, the extremal representative satisfies
\[
\deg_{L(\mathcal F)}(f)=k(D-1)=kD-k. \tag{1}
\]
By \noderef{KUniformKSimplePairLowerBound},
\[
P_{L(\mathcal F)}(f)\ge \frac{k-1}{2}D^2-k^2D. \tag{2}
\]
Set
\[
a=\frac{k-1}{2},\qquad
y=\frac{P_{L(\mathcal F)}(f)}{D^2}.
\]
Then (2) says
\[
y\ge a-\frac{k^2}{D}. \tag{3}
\]
Increasing $D_0$ once more, we may and do assume
$k^2/D\le a$, so the lower endpoint in (3) is nonnegative.

Independent triples in $N_{L(\mathcal F)}(f)$ are triangles in the complement
of the graph induced by $L(\mathcal F)$ on that neighborhood.  The complement
is $k$-partite.  Indeed, \noderef{KUniformKSimpleExtremalSaturation} also
says that each of the $k$ vertices of $f$ has degree $D$, so there are
$k(D-1)$ incidences between vertices of $f$ and edges other than $f$.  The
same node says that there are exactly $k(D-1)$ distinct neighboring edge
indices.  Since every neighbor contains at least one vertex of $f$, each
neighbor contains exactly one such vertex.  Partition the neighbors according
to that unique vertex.  Two neighbors in the same part share that vertex, so
they are adjacent in the line graph and therefore nonadjacent in the
complement; hence the complement is $k$-partite.  The number of edges of this
complement is precisely $P_{L(\mathcal F)}(f)$ by
\noderef{IndependentPairCount}.  Hence
\noderef{KPartiteTriangleCount} gives
\[
T_{L(\mathcal F)}(f)
\le
\binom{k}{3}
\left(\frac{P_{L(\mathcal F)}(f)}{\binom{k}{2}}\right)^{3/2}. \tag{4}
\]

Also,
\[
P_{L(\mathcal F)}(f)\le \binom{\deg_{L(\mathcal F)}(f)}2
=\binom{kD-k}{2}<\frac{k^2D^2}{2},
\]
because $P$ counts independent pairs among the neighbors of $f$.  Therefore
\[
0\le y<\frac{k^2}{2}. \tag{5}
\]

Substitute (1), the definition of $y$, and (4) into
\noderef{LocalBParameter}:
\[
\begin{aligned}
b_{L(\mathcal F),kD}(f)
&\le
\frac{kD-k}{kD}
-\frac{P_{L(\mathcal F)}(f)}{(kD)^2}\\
&\quad+
\frac{\binom{k}{3}
\left(P_{L(\mathcal F)}(f)/\binom{k}{2}\right)^{3/2}}
{(kD)^3}.
\end{aligned}
\]
Since $P_{L(\mathcal F)}(f)=yD^2$, this becomes
\[
\begin{aligned}
b_{L(\mathcal F),kD}(f)
&\le
1-\frac1D-\frac{y}{k^2}
+c_k y^{3/2}, \tag{6}
\end{aligned}
\]
where
\[
c_k=\frac{\binom{k}{3}}{k^3\binom{k}{2}^{3/2}} .
\]

It remains only to check the one-variable estimate on the whole feasible
interval.  Define
\[
F(x)=1-\frac{x}{k^2}+c_k x^{3/2}
\qquad (x\ge0).
\]
For $0\le x\le k^2/2$,
\[
F'(x)=-\frac1{k^2}+\frac32 c_k\sqrt{x}
\le
-\frac1{k^2}+\frac32 c_k\frac{k}{\sqrt2}.
\]
A direct simplification gives
\[
\frac32 c_k\frac{k}{\sqrt2}
=\frac{k-2}{2k^{5/2}\sqrt{k-1}}
\le \frac1{k^2},
\]
because $k-2\le 2\sqrt{k(k-1)}$ for every $k\ge2$.  Hence $F$ is
nonincreasing on $[0,k^2/2]$.

By (3) and (5), the value $y$ in (6) lies in the interval
$[a-k^2/D,k^2/2]$.  Since $F$ is nonincreasing there,
\[
F(y)\le F\left(a-\frac{k^2}{D}\right).
\]
On the interval $[a-k^2/D,a]$ the same derivative formula gives
$F'(x)\ge -1/k^2$, and so
\[
F\left(a-\frac{k^2}{D}\right)
\le F(a)+\frac1D.
\]
Combining this with (6) yields
\[
b_{L(\mathcal F),kD}(f)\le F(a).
\]
Finally,
\[
\begin{aligned}
F(a)
&=
1-\frac{k-1}{2k^2}
+\frac{\binom{k}{3}}{k^3\binom{k}{2}^{3/2}}
\left(\frac{k-1}{2}\right)^{3/2}\\
&=
1-\frac{k-1}{2k^2}
+\frac{(k-1)(k-2)}{6k^3\sqrt{k}},
\end{aligned}
\]
where the last equality is the cancellation
\[
\frac{\binom{k}{3}\left((k-1)/2\right)^{3/2}}
{k^3\left(k(k-1)/2\right)^{3/2}}
=\frac{(k-1)(k-2)}{6k^3\sqrt{k}}.
\]
Thus the desired bound holds for the extremal representative with room to
spare, and therefore also with the extra $+\eta$ in the statement.  The
finite local-pattern reduction at the start transfers the same bound back to
the original pair $(\mathcal H,e)$.
\end{proof}
